\section{The statements of Part I and a conserved current}\label{sec:current}

Part I reproduces the proof by OpenAI \cite{OAI26} of Sogge's local smoothing conjecture and
of its consequence for Bochner--Riesz means. Throughout Part I the spatial dimension is
$n\ge2$, we put $d=n-1$, and
\begin{gather*}
U(T)=e^{-iT|D|},\qquad I=[1,2],\qquad J=(1-\Delta)^{1/2},\\
s_+(p)=\max\Bigl\{0,\frac{n-1}{2}-\frac np\Bigr\},\qquad \pc=\frac{2n}{n-1}.
\end{gather*}

\begin{theorem}[Local smoothing \cite{OAI26}]\label{thm:ls}
For $n\ge2$, $2\le p<\infty$ and $\eps>0$,
\[
\|U(T)f\|_{L^p(\R^n\times I)}\le C_{n,p,\eps}\|J^{s_+(p)+\eps}f\|_{L^p(\R^n)} .
\]
\end{theorem}

The whole proof rests on one estimate at the exponent $\pc$, on a single frequency annulus.

\begin{theorem}[Critical annular estimate \cite{OAI26}]\label{thm:annular}
Let $n\ge2$, $\varphi\in C^\infty_c(\R^n\setminus\{0\})$ and $P_N=\varphi(D/N)$. For $N\ge2$,
\[
\|U(T)P_Nf\|_{L^{\pc}(\R^n\times I)}\le C_{n,\varphi}\,[\log(2+N)]^{2/\pc}\,\|f\|_{L^{\pc}(\R^n)} .
\]
\end{theorem}

\subsection{The plan}\label{ssec:plan}
Put $h=N^{-1}$ and $a=Kh$ with a large fixed $K$. Besides the time $T$ there will be two
auxiliary real time variables $\tau\in\R^2$, which come from an imaginary displacement of
the wave along light rays. For data $f$ on a fixed torus whose frequencies lie in one small
cap of the annulus $|\xi|\asymp h^{-1}$ and with $\|f\|_2=1$, we build a nonnegative energy
$\mathcal F(T,\tau)$ and prove, with $\mathcal H(T,s)$ the mean of $\mathcal F$ over
$|\tau|=s$,
\[
\int_1^2\|U(T)f\|_{\pc}^{\pc}\,dT\lesssim a^{-2}\int_1^2\mathcal H(T,2a)\,dT
\lesssim[\log(2/h)]^2a^{-2}\mathcal H(0,4a)\lesssim[\log(2/h)]^2\|f\|_{\pc}^{\pc}.
\]
The outer inequalities compare the energy with the wave (Section \ref{sec:fiber}); the middle
one uses only the differential inequality $\Delta_{T,\tau}\mathcal F\le Cs^{-2}e^{-cs/h}\mathcal F$
on $2a\le|\tau|\le3$ (Section \ref{sec:cylinder}).

The energy is built as follows. A finite sum of plane waves whose directions lie in one cap
is continued to complex spacetime. Translation by $-i\ell(1,\omega)$ along a light ray
multiplies a mode of frequency $\xi$ by $\exp(-\ell|\xi||\omega-\xi/|\xi||^2/2)$, a Gaussian
of angular width $(N\ell)^{-1/2}$. The fluxes of the symmetric first-order form of the wave
equation produce such displacements, and two flows built from the continued wave give a
positive semidefinite tensor with divergence-free columns whose time component still sees
the wave on the cap (this section). On each spatial fibre, the time component $\mu$ of a
regularised version of this tensor is used as the density of a quaternionic Monge--Amp\`ere
equation $\det\Hfib(\lambda|\Phi|^2+u)=\lambda\mu$ with zero boundary values, where $\Phi$
is a complex scalar constant along the current and $\lambda$ is large. The background
$\lambda|\Phi|^2$ has a Hessian of rank one; adding $\alpha I$ to it changes its determinant to
$(\beta+\alpha)\alpha^{n-1}$ with $\beta\asymp\lambda$. Hence a density $\mu\ge m$ on a ball
forces the depth of $u$ at its centre to be at least a multiple of the squared radius times
$m^{1/(n-1)}$, and $\int(-u)\mu$ detects $m^{n/(n-1)}$. Since $\mu$ is a square of the wave,
the exponent is $2n/(n-1)=\pc$. A correction of Aubin--Mabuchi type makes the full determinant
cancel in $\Delta_{T,\tau}$ of the energy, and the only positive contribution left is the
wave flux through the side of the fibre, which carries a Gaussian factor $e^{-cs/h}$.

\subsection{Null coordinates and modes}\label{ssec:modes}
Choose an orthonormal frame $X=(x,b)\in\R^d\times\R$ of $\R^n$ and put
\[
t=(T+b)/\sqrt2,\qquad s_0=(b-T)/\sqrt2 .
\]
In these coordinates the wave operator $\partial_T^2-\Delta_X$ becomes, up to sign,
$2\partial_t\partial_{s_0}+\Delta_x$. A plane wave with phase $X\cdot\xi-T|\xi|$ and
$\xi_b+|\xi|>0$ is
\begin{equation}\label{eq:mode}
u_{\kappa,v}=c_{\kappa,v}\exp\Bigl(i\kappa\bigl(s_0+v\cdot x-\tfrac12|v|^2t\bigr)\Bigr),\qquad
\kappa=\frac{\xi_b+|\xi|}{\sqrt2},\qquad v=\frac{\xi_x}{\kappa}.
\end{equation}
Indeed, the phase is $x\cdot\xi_x+s_0(\xi_b+|\xi|)/\sqrt2+t(\xi_b-|\xi|)/\sqrt2$, and
$(\xi_b-|\xi|)/\sqrt2=-|\xi_x|^2/(2\kappa)=-\kappa|v|^2/2$. Throughout the construction
$u_{\rm w}=\sum_{\kappa,v}u_{\kappa,v}$ is a finite sum with $\kappa>0$ and complex coefficients.

Index the states by $0,\dots,d$, let $E_{ab}$ be the elementary $n\times n$ matrices, and
in the real spacetime coordinates $\mathsf X=(t,x,s_0)$ put
\[
C^t=2E_{00},\qquad C^{x_j}=E_{0j}+E_{j0},\qquad C^{s_0}=-\sum_{j=1}^dE_{jj},\qquad
\mathscr P=\sum_\mu C^\mu\partial_{\mathsf X_\mu}.
\]
For a scalar $\psi$ the vector $(\partial_{s_0}\psi,\nabla_x\psi)$ satisfies
$\mathscr P(\partial_{s_0}\psi,\nabla_x\psi)^{\rm t}=((2\partial_t\partial_{s_0}+\Delta_x)\psi,0,\dots,0)^{\rm t}$:
the first entry is $2\partial_t\partial_{s_0}\psi+\sum_j\partial_{x_j}\partial_{x_j}\psi$, and
the $j$th entry is $\partial_{x_j}\partial_{s_0}\psi-\partial_{s_0}\partial_{x_j}\psi=0$. At the
covector $p=\kappa(-|v|^2/2,v,1)$ of \eqref{eq:mode} the symbol $C(p)=\sum_\mu p_\mu C^\mu$
satisfies
\begin{equation}\label{eq:symbol}
-C(p)=\kappa\begin{pmatrix}|v|^2&-v^{\rm t}\\-v&I_d\end{pmatrix},\qquad
-w^{\rm t}C(p)w=\kappa|w'-vw_0|^2\quad(w=(w_0,w')\in\R^n).
\end{equation}
Its kernel is spanned by $E(v)=(1,v)^{\rm t}$, and
$f_{\rm sys}=\sum_{\kappa,v}\kappa^{d/2}E(v)u_{\kappa,v}$ solves $\mathscr Pf_{\rm sys}=0$.

\subsection{The Gaussian from an imaginary displacement}
Translating $(T,X)$ by $-i\varrho(1,\omega)$ with $\varrho>0$ and $|\omega|=1$ multiplies a mode of
frequency $\xi$ by
\[
e^{\varrho(\xi\cdot\omega-|\xi|)}=e^{-\varrho|\xi||\hat\xi-\omega|^2/2},\qquad\hat\xi=\xi/|\xi| ,
\]
since $|\hat\xi-\omega|^2=2-2\hat\xi\cdot\omega$. The system produces such light rays as quadratic
fluxes. For a real state $w=(w_0,w')$ let $q(w)^\mu=w^{\rm t}C^\mu w/2$; in null coordinates
$q(w)=(w_0^2,w_0w',-|w'|^2/2)$. Since $\partial_t=(\partial_T+\partial_b)/\sqrt2$ and
$\partial_{s_0}=(\partial_b-\partial_T)/\sqrt2$, its physical components satisfy
$q(w)^T=|q(w)^X|=(w_0^2+|w'|^2/2)/\sqrt2$, so $q(w)$ is a light ray; at $w=E(v)$ it is
$(|\xi|,\xi)/\kappa$. We use the sum of the fluxes of two real states,
\begin{equation}\label{eq:twoflux}
q^\mu=q(w_1)^\mu+q(w_2)^\mu=\frac12\sum_{\ell=1}^2w_\ell^{\rm t}C^\mu w_\ell,\qquad w_1,w_2\in\R^n .
\end{equation}
At $\mathsf X-iq$ a mode acquires the factor $e^{p\cdot q}$. Introduce a second imaginary unit
$\kk$, with $\kk^2=-1$, put $\Ck=\{\alpha+\kk\beta:\alpha,\beta\in\R\}$, and pair the two states
into
\[
z=w_1+\kk w_2=(R,\zeta)\in\Ck\times\Ck^d .
\]
Then $q^\mu=z^*C^\mu z/2$, where $^*$ is the conjugate transpose, and adding the two real
squares of \eqref{eq:symbol} gives
\begin{equation}\label{eq:gauss}
e^{p\cdot q}=\exp\Bigl(-\frac\kappa2|\zeta-vR|^2\Bigr).
\end{equation}
For each $R$ the square of $\kappa^{d/2}e^{p\cdot q}$ has integral $\pi^d$ in $\zeta$, which
explains the factor $\kappa^{d/2}$ in $f_{\rm sys}$.

\subsection{Two auxiliary flows}
For constant real symmetric matrices $C^\mu$ the bilinear flux of two states is
$j^\mu(a,b)=a^{\rm t}C^\mu b$, and its divergence for two fields is
$(\mathscr Pa)^{\rm t}b+a^{\rm t}\mathscr Pb$; this is the conservation law of symmetric
hyperbolic systems \cite{Fri54}. The derivative of $q^\mu$ in the direction $\delta w_\ell$ is
$j^\mu(w_\ell,\delta w_\ell)$, and this suggests the following vector fields.

\begin{lemma}[\cite{OAI26}]\label{lem:flows}
Let $\mathscr P=\sum_\mu C^\mu\partial_{\mathsf X_\mu}$ have constant real symmetric
coefficients, let $f_{\rm sys}$ be an entire solution of $\mathscr Pf_{\rm sys}=0$, and let $q$ be
the displacement \eqref{eq:twoflux}. In the real coordinates $(\mathsf X,w_1,w_2)$ put
\[
G=\overline{f_{\rm sys}(\mathsf X-iq)},\quad V_{\ell a}=\sum_\mu(C^\mu w_\ell)_a\,\partial_{\mathsf X_\mu},\quad
D_{\ell a}=\partial_{w_{\ell a}}+iV_{\ell a},\quad \mathcal Z_\ell=\sum_aG_aD_{\ell a}.
\]
Then the real tensor
\begin{equation}\label{eq:B}
B=\Re\sum_{\ell=1}^2\mathcal Z_\ell\otimes\overline{\mathcal Z_\ell}
\end{equation}
is symmetric and positive semidefinite, and each of its columns has zero divergence.
\end{lemma}

\begin{proof}
The field $D_{\ell a}$ annihilates every component of $\mathsf X-iq$: it sends
$\mathsf X_\mu-iq^\mu$ to $-i(C^\mu w_\ell)_a+i(C^\mu w_\ell)_a=0$. Since $f_{\rm sys}$ is
holomorphic, $D_{\ell a}\overline G=0$, that is
\[
\partial_{w_{\ell a}}G=iV_{\ell a}G .
\]
The coefficients of $V_{\ell a}$ depend only on $w$, so $D_{\ell a}$ has zero divergence, and
\[
\operatorname{div}\mathcal Z_\ell=\sum_aD_{\ell a}G_a=\sum_a(\partial_{w_{\ell a}}G_a+iV_{\ell a}G_a)
=2i\sum_aV_{\ell a}G_a=2i\,w_\ell^{\rm t}\,\overline{(\mathscr Pf_{\rm sys})(\mathsf X-iq)}=0 .
\]
For the flat connection the identity $D_{\ell a}\overline G=0$ shows that, in
$\nabla_{\mathcal Z_\ell}\overline{\mathcal Z_m}=\sum_{a,b}G_a D_{\ell a}(\overline{G_b})\,\overline{D_{mb}}
+\sum_{a,b}G_a\overline{G_b}\,D_{\ell a}(\overline{D_{mb}})$, only the derivatives of the
coefficients of $\overline{D_{mb}}=\partial_{w_{mb}}-iV_{mb}$ remain. These are
$-i\,\delta_{\ell m}\sum_\mu (C^\mu)_{ba}\partial_{\mathsf X_\mu}$, so
\[
\nabla_{\mathcal Z_\ell}\overline{\mathcal Z_m}=-i\,\delta_{\ell m}\sum_\mu\bigl(G^{\rm t}C^\mu\overline G\bigr)\partial_{\mathsf X_\mu}.
\]
The vector on the right is real since the $C^\mu$ are real symmetric, so its product with $-i$ is
purely imaginary. The column divergence of $\mathcal Z_\ell\otimes\overline{\mathcal Z_\ell}$ is
$(\operatorname{div}\mathcal Z_\ell)\overline{\mathcal Z_\ell}+\nabla_{\mathcal Z_\ell}\overline{\mathcal Z_\ell}$,
whose real part vanishes. Finally, if $\mathcal Z_\ell=A_\ell+iB_\ell$ with real fields, its
contribution to \eqref{eq:B} is $A_\ell\otimes A_\ell+B_\ell\otimes B_\ell\ge0$.
\end{proof}

The imaginary part of the last identity at $m=\ell$ gives
$[A_\ell,B_\ell]=\sum_\mu(G^{\rm t}C^\mu\overline G)\partial_{\mathsf X_\mu}$, which identifies the
physical current inside the lifted variables. We now express the four real squares of $B$ in
coordinates attached to spacetime.

\subsection{Quaternionic coordinates adapted to the wave}\label{ssec:quat}
The quaternions $\HH$ have units $\ii,\jj,\kk$ with $\ii^2=\jj^2=\kk^2=-1$, $\ii\jj=\kk$ and
$\kk\ii=\jj$. The Fourier complex plane is identified with $\Ci=\R+\R\ii$; every quaternion is
uniquely $\alpha+\kk\beta$ with $\alpha,\beta\in\Ci$, and the planes $\Ci$ and $\kk\Ci$ are
orthogonal for $\ip uv=\Re(u\bar v)$. Columns carry scalars acting on the right, and $^*$
denotes the quaternionic conjugate transpose. We look for one auxiliary pair for each of the
$n+1$ real spacetime coordinates and write each triple as $\mathsf X_\mu+\omega_\mu\ii$ with
$\omega_\mu\in\Ck$; these triples lie in $\mathcal D=\operatorname{span}_\R(1,\ii,\jj)$.

For a quaternionic Hermitian matrix $A$ on spacetime indices put
\begin{equation}\label{eq:AR}
A^{\R}_{\mu\alpha,\nu\beta}=\Re(\overline{e_\alpha}A_{\mu\nu}e_\beta),\qquad(e_0,e_1,e_2)=(1,\ii,\jj).
\end{equation}
If $A=\chi\chi^*$ this real tensor is the Gram matrix of the rows $\overline{e_\alpha}\chi_\mu$,
so the three rows over each $\mu$ are $Q,-\ii Q,-\jj Q$.

In source coordinates write the four squares of \eqref{eq:B} as a Gram matrix by setting
$Q_I=\mathcal Z_1^I+\kk\mathcal Z_2^I$ for each coordinate $I$. Orthogonality of $\Ci$ and
$\kk\Ci$ gives $B^{IJ}=\Re(Q_I\overline{Q_J})$. With $M_{\mu a}=(C^\mu z)_a$ the rows in the
real spacetime and in the state directions are
\[
Q_{\mathsf X_\mu}=\sum_aM_{\mu a}\,\ii\, G_a,\qquad(Q_{w_{1a}},Q_{w_{2a}})=(G_a,\kk G_a).
\]
For a holomorphic coordinate $\omega_\mu(z)$ let $L_{\mu a}=\partial_{z_a}\omega_\mu$.
Multiplication by $c=c_1+\kk c_2$ on an auxiliary pair has real rows $(c_1,-c_2)$ and $(c_2,c_1)$,
so it sends $(G_a,\kk G_a)$ to $(\bar cG_a,\kk\bar cG_a)$; the new pair over $\mu$ therefore
starts with $\sum_a\overline{L_{\mu a}}G_a$. The choice $L=M$ gives exactly $-\ii Q_{\mathsf X_\mu}$,
because $-\ii c\,\ii=\bar c$ for $c\in\Ck$, and its partner is $\kk(-\ii Q_{\mathsf X_\mu})=-\jj Q_{\mathsf X_\mu}$.

\subsubsection*{The holomorphic flux map}
By symmetry, $L_{\mu a}=(C^\mu z)_a$ integrates to the quadratic $z^{\rm t}C^\mu z/2$: the
Hermitian square $z^*C^\mu z/2$ weights the modes, and the bilinear square moves their Gram
rows. There are $n$ complex state coordinates and $n+1$ spacetime coordinates, so add a
passive variable $\eta\in\Ck$ in whose two real directions the flows have no component. For
a constant real symmetric system as in Lemma \ref{lem:flows} and a fixed $c\in\R^{n+1}$ set
\[
\Theta:(z,\eta)\longmapsto\omega,\qquad\omega_\mu=\tfrac12z^{\rm t}C^\mu z+c^\mu\eta .
\]
It fixes the real spacetime coordinates, and on a local inverse chart its real volume Jacobian
is $J=|\det_{\C}[M,c]|^2$. Push $B$ forward as a density,
\begin{equation}\label{eq:pushforward}
\widetilde B\circ\Theta=J^{-1}\,D\Theta\,B\,D\Theta^{\rm t}.
\end{equation}
The differentials transform the Gram rows as described, and $J^{-1}$ converts source volume to
target volume, so
\begin{equation}\label{eq:chi-tilde}
\widetilde B=(\widetilde\chi\widetilde\chi^*)^{\R},\qquad
\widetilde\chi_\mu\circ\Theta=J^{-1/2}\sum_aM_{\mu a}\ii G_a .
\end{equation}
As a function of $\omega$, the passive coordinate satisfies $(\partial_\omega\eta)M=0$, because
varying $z$ at fixed $\eta$ produces the columns of $M$; hence $(\partial_\omega\eta)\widetilde\chi=0$.

The tensor $\widetilde B$ is again conserved. Test its $K$th column against the gradient of a
real $\varphi$ with compact support in the target chart. Changing variables in
\eqref{eq:pushforward}, the weak divergence becomes
$-\sum_{I,L}\int B^{IL}(\partial_L\Theta^K)\partial_I(\varphi\circ\Theta)$, and integration by
parts in the source gives
\[
J(\operatorname{div}\widetilde B)^K\circ\Theta=\sum_L(\partial_L\Theta^K)(\operatorname{div}B)^L
+\sum_{I,L}B^{IL}\partial_I\partial_L\Theta^K .
\]
The first sum vanishes by Lemma \ref{lem:flows}. In the second, the two diagonal state blocks
of $B$ both equal $(\Re(G_a\overline{G_b}))_{ab}$ and the mixed state block vanishes, while each
real component of the holomorphic map satisfies
$\partial_{w_{1a}}\partial_{w_{1b}}\Theta^K+\partial_{w_{2a}}\partial_{w_{2b}}\Theta^K=0$. The other
coordinates enter $\Theta$ affinely, so the second sum vanishes too.

\subsubsection*{The wave in these coordinates}
For the wave matrices add $\eta$ in the $s_0$ coordinate. In the null frame the flux map is
\begin{equation}\label{eq:fluxmap}
\Theta:(R,\zeta,\eta)\longmapsto(Y,P,S_1)=(R^2,R\zeta,\eta-\zeta\cdot\zeta/2).
\end{equation}
If $(\tau,P,Z)$ are the target auxiliary coordinates in the physical frame $(T,x,b)$, then
$Y=(\tau+Z)/\sqrt2$ and $S_1=(Z-\tau)/\sqrt2$. On $Y\ne0$, inverting \eqref{eq:fluxmap} gives
$R=\sqrt Y$ on a local branch and $\zeta=P/R$, and the single-valued coordinates are
\begin{equation}\label{eq:WPhi}
W=\frac PY=\frac\zeta R,\qquad\Phi=S_1+\frac{P\cdot P}{2Y}=\eta,\qquad y=|Y| .
\end{equation}
All dots between complex vectors are bilinear, and $|\cdot|$ is the Euclidean norm. The complex
Jacobian determinant of $\Theta$ is $2R^{d+1}$, so
\begin{equation}\label{eq:jacobian}
J=4|R|^{2d+2}.
\end{equation}
Under $M$ a polarisation vector $E(v)$ maps to
\begin{equation}\label{eq:Vv}
ME(v)=(2R,\zeta+Rv,-\zeta\cdot v)^{\rm t}=2RV^v,\qquad V^v=\bigl(1,(v+W)/2,-W\cdot v/2\bigr)^{\rm t}.
\end{equation}
Put $\ell(w)=(-w\cdot w/2,w,1)$, so that $\kappa\ell(v)$ is the real phase covector; one checks
$\ell(w)V^v=-(w-W)\cdot(w-v)/2$. The column $V^v$ thus lies in the kernels of both $\ell(v)$ and
$\partial\Phi=\ell(W)$, and differentiating $W=P/Y$ gives the second identity in
\begin{equation}\label{eq:tangency}
(\partial\Phi)V^v=0,\qquad(\partial W)V^v=\frac{v-W}{2Y} .
\end{equation}
These are complex differentials in $(Y,P,S_1)$. The current stays on a level set of $\Phi$,
and each mode moves $W$ in proportion to its angular mismatch $v-W$.

For $e=Y/y$ put $\iota_e(\alpha+i\beta)=\alpha+e\ii\beta$. This embeds the Fourier complex
numbers isometrically into $\HH$, since $\ii c=\bar c\ii$ for $c\in\Ck$ gives $(e\ii)^2=-1$.
Define the positive weight
\[
m_{\kappa,v}=(\kappa/y)^{d/2}\exp\Bigl(-\frac{\kappa y}2|W-v|^2\Bigr);
\]
the factor $y^{-d/2}$ comes from the change of volume.

\begin{proposition}[The wave current \cite{OAI26}]\label{prop:current}
For a finite superposition of modes \eqref{eq:mode} set
\begin{equation}\label{eq:chi}
\chi=\sum_{\kappa,v}V^v\,\iota_e\bigl(\overline{u_{\kappa,v}}\bigr)\,m_{\kappa,v},\qquad A^0=\chi\chi^* .
\end{equation}
On $Y\ne0$ the matrix $A^0$ is smooth and positive semidefinite, $(\partial\Phi)A^0=0$, and each
column of $(A^0)^{\R}$ has zero real divergence. The same holds after the orthogonal change
from null to physical coordinates, and
\[
V^v_T=(1+W\cdot v/2)/\sqrt2,\qquad A^0_{TT}=|\chi_T|^2 .
\]
\end{proposition}

\begin{proof}
We identify the column \eqref{eq:chi-tilde} on one inverse chart. By \eqref{eq:gauss} and
\eqref{eq:Vv} the contribution of a single mode to $\widetilde\chi$ is
$V^v g\,\ii\,\overline{u_{\kappa,v}}\,m_{\kappa,v}$ with $g=R/|R|$, since
$|\zeta-vR|^2=y|W-v|^2$ and $2|R|/\sqrt J=y^{-d/2}$. Multiplying every row on the right by the
unit quaternion $(g\ii)^{-1}$ preserves the Gram matrix; conjugation by $g$ takes $\ii$ to
$g^2\ii=e\ii$, and $\ii$ commutes with the Fourier amplitudes, so the column becomes
\eqref{eq:chi}. This formula involves $Y$ itself, hence is smooth and independent of the
branch of the square root. Positivity follows from the outer product form, tangency from
\eqref{eq:tangency}, and conservation from the discussion of $\widetilde B$. The change to
physical coordinates acts by the same real orthogonal matrix in each of the three coordinate
slots, so it preserves column divergence and the pairing of rows and columns. Finally
$T=(t-s_0)/\sqrt2$ gives $V^v_T$.
\end{proof}

For a single $\kappa$ and real $W$ the exponential in $m_{\kappa,v}$ is a Gaussian window in the
transverse frequency $\kappa v$ centred at $\kappa W$, and since $V^v_t=1$ the null component
$A^0_{tt}$ is, up to normalisation, the Husimi density of the scalar wave
\cite{Hus40,KLO16}, a phase-space density of the kind also used for quantum states on
symplectic manifolds \cite{Bha26a}. For the full wave, \eqref{eq:chi} sums the amplitudes
over $\kappa$ before taking the square, which keeps the interference that produces a focus.

\subsection{The time component on a frequency cap}\label{ssec:cap}
At a fixed auxiliary point write $V^v_T=b_0+\kk b_1$, so that
\[
b_0=\frac{1+\Re(W)\cdot v/2}{\sqrt2},\qquad b_1=\frac{\Imk(W)\cdot v}{2\sqrt2},
\qquad\Imk(\alpha+\kk\beta)=\beta .
\]
Let $S_\ell$ be the scalar spatial multiplier sending each mode of \eqref{eq:mode} to
$b_\ell m_{\kappa,v}u_{\kappa,v}$ at this auxiliary point. By \eqref{eq:chi},
\begin{equation}\label{eq:ATT}
\chi_T=\iota_e\bigl(\overline{S_0u_{\rm w}}\bigr)+\kk\,\iota_e\bigl(\overline{S_1u_{\rm w}}\bigr),\qquad
A^0_{TT}=|S_0u_{\rm w}|^2+|S_1u_{\rm w}|^2 ,
\end{equation}
because the two terms lie in the conjugates of the orthogonal planes $\Ci$ and $\kk\Ci$ by a
square root of $e$ in $\Ck$. At $W=0$ the second multiplier vanishes and $2A^0_{TT}$ is exactly
the squared modulus of the wave with weight $(\kappa/y)^{d/2}\exp(-\kappa y|v|^2/2)$ on each mode.

For the estimates take $u_{\rm w}=e^{-iT|D|}f$ on the fixed torus $\T^n=\R^n/(100\Z)^n$, with the
axis $b$ along the centre of a frequency cap; the torus itself is not rotated. For $0<h\le1$
suppose $\hat f$ is supported in
\[
c_0/h\le|\xi|\le C_0/h,\qquad|v|\le\gamma,
\]
with fixed $0<c_0<C_0$ and a small aperture $\gamma>0$ independent of $h$. The identity
$|\xi|=\kappa(1+|v|^2/2)/\sqrt2$ gives $\kappa\asymp h^{-1}$. At $y\asymp h$ and $|W|,|v|\le1/10$
the Gaussian and $b_0$ are bounded below, so $h^dS_0$ has a symbol bounded above and below on
the cap; for larger $y/h$ the cost of inverting $S_0$ depends on the mismatch $|W-v|$. There are
$O(h^{-n})$ frequencies in the annulus, uniformly in the frame, so Cauchy--Schwarz and Parseval
give $\sum_{\kappa,v}|c_{\kappa,v}|\lesssim h^{-n/2}\|f\|_2$. The formulas in null coordinates may
be used on the covering space; the resulting current is periodic in the original spatial
variables. Nonzero cap data are normalised by $\|f\|_2=1$ in the energy estimates.
